\section{Related Work}
\label{sec:related_work}

\subsection{UMI Deduplication Tools}

UMI deduplication tools group raw reads by UMI identity and produce a single consensus read per molecule group. UMI-tools \autocite{Smith17umitools} introduced network-based grouping methods (directional, adjacency, and cluster) where nodes represent UMIs and edges connect UMIs within a fixed edit distance. This corrects for sequencing errors in the UMI itself and remains the most widely cited method in the literature. HUMID \autocite{Langedijk22humid} takes a reference-free approach, using a trie over UMI sequences to perform Hamming-distance grouping directly on FASTQ input without prior alignment. fgbio \autocite{fgbio} is the de facto clinical standard for duplex UMI workflows: its \texttt{GroupReadsByUmi} and \texttt{CallDuplexConsensusReads} stages produce a single-strand consensus per strand and a duplex consensus where both strands agree. UMICollapse \autocite{Liu19umicollapse} achieves order-of-magnitude speed improvements over UMI-tools using N-gram indexing and BK-trees. Commercial tools such as the Sentieon ctDNA pipeline offer a 2--3$\times$ end-to-end speedup at the cost of being closed source.

All of these tools share a common output model: a deduplicated FASTQ or BAM file where each record represents one consensus read. The molecule identity that produced each consensus, including whether it was duplex or simplex, how many raw reads supported it, and which strand each read came from, is either discarded entirely or encoded in SAM auxiliary tags that downstream tools rarely use. In fgbio's own benchmarks, consensus base quality scores generated by \texttt{CallDuplexConsensusReads} were not propagated into variant calling \autocite{fgbio}, which means the duplex-level confidence information was computed and then ignored.

kam differs from this family by making the molecule the primary unit of analysis throughout the pipeline. Molecule provenance, duplex status, strand family size, and per-base confidence are not metadata appended to a read record; they are first-class fields in the \texttt{MoleculeEvidence} type that flows from assembly through k-mer indexing to variant calling. This propagation is what enables duplex-aware scoring at the calling stage.

\subsection{K-mer-Based Variant Detection}

The km/kmtools pipeline \autocite{Bourgey19km} uses Jellyfish \autocite{Marcais11jellyfish} k-mer counts as the basis for alignment-free variant detection. Jellyfish counts k-mer occurrences in a set of sequences using a lock-free hash table and outputs per-k-mer integer counts. km walks de Bruijn graph paths built from those counts to find sequence variants: it identifies k-mers present in a sample but absent or underrepresented in a reference, then assembles those k-mers into candidate variant sequences. This approach handles insertions, deletions, and SNVs without reference alignment, avoiding alignment bias in repetitive regions. kmtools parallelises km across a target panel and aggregates results.

The core limitation is that Jellyfish stores integer counts only. It has no concept of read origin, molecule identity, or quality. The canonical k-mer flag (\texttt{-C}) collapses both strands into the same counter, erasing strand information entirely. A k-mer count of 10 could represent ten PCR duplicates from a single molecule or ten independent molecules, and Jellyfish cannot distinguish them. Quality scores are ignored, so a k-mer supported by ten Q10 reads looks identical to one supported by ten Q40 reads. km's variant calling uses a global ratio threshold (\texttt{--ratio}) with no statistical model: the same threshold is applied regardless of sequencing depth, which means the false positive rate changes with coverage even when the true VAF is constant.

kam differs from this family by replacing integer k-mer counts with \texttt{MoleculeEvidence}: each k-mer entry records the number of distinct molecules that contributed it, broken down by duplex, simplex-forward, and simplex-reverse. Variant calling operates on molecule counts under a binomial model with a per-site background, not on a global ratio threshold. Duplex molecules, whose error rate is approximately $10^{-7}$ compared to $10^{-2}$ for a single read, receive higher statistical weight. This allows depth-calibrated confidence intervals and strand-bias detection as a first-class filter.

\subsection{Alignment-Based Low-VAF Variant Callers}

The dominant approach to somatic variant detection in ctDNA is to align reads to a reference genome and call variants from the resulting BAM file. GATK Mutect2 \autocite{McKenna10gatk} uses a haplotype assembly model with a somatic-specific likelihood that accounts for the possibility that a variant is real but rare. VarDict \autocite{Lai16vardict} detects variants at low frequency using local realignment and a Fisher's exact test for strand bias. LoFreq \autocite{Wilm12lofreq} models per-base sequencing error rates explicitly and uses a Poisson binomial model to call variants that would be indistinguishable from noise in a standard caller. DeepVariant \autocite{Poplin18deepvariant} reformulates variant calling as an image classification problem, encoding pileup evidence as a tensor and using a convolutional neural network to produce genotype posteriors.

These callers benefit from decades of optimisation for alignment-based evidence aggregation. However, alignment imposes limitations that matter for ctDNA at 0.1\% VAF. Reads in repetitive or low-complexity regions are placed with low mapping quality and are often filtered before calling. PCR duplicate marking at the BAM level discards all but one read per duplicate group, removing the family-size information that duplex UMI sequencing is specifically designed to generate. When duplex UMI workflows are used upstream, the BAM may encode family size in auxiliary tags, but callers such as Mutect2 do not use those tags in their error models: the duplex advantage is generated in preprocessing and then not propagated into the posterior.

kam differs from this family by operating without alignment. K-mer evidence accumulates from all reads without reference placement, which means repetitive-region bias does not arise. More directly, because kam tracks molecules rather than reads, there is no separate duplicate-marking step and no information loss: the molecule count used in the binomial likelihood is the count of distinct original DNA molecules, not the count of reads after duplicate removal.
